\begingroup\fontsize{10}{12.1}\selectfont
\setlength{\parskip}{1.5pt}\setlength{\abovedisplayskip}{4pt}\setlength{\belowdisplayskip}{4pt}
\subsection{Angular coordinates, winding normalization, and action units}
\label{sec:ii-angulos}

The action quantity and the angular coordinate belong to different
spaces. The section \(\hbar_{\rm ret}\) occupies the action line
\(\mathcal L_S\); \(\eta_{\rm ret}\) is its coefficient in the
declared basis. Angular representation acts on the dimensionless pair
\((\alpha,\eta_{\rm ret})\), after its construction:
\begin{equation}
 \mathcal P(\alpha,\eta_{\rm ret})
 =\bigl([1000\alpha]_{360},[2(\eta_{\rm ret}+\alpha)]_{360}\bigr).
 \label{eq:ii-par-angular}
\end{equation}
Nonadic completion determines the factor 1000; bidirectional
closure determines the factor two. The circular chart of 360
units makes the partitions
\(12\cdot30=4\cdot90=3\cdot120\) visible. On the positive branch of
the article, the representatives are
\[
 A_{\deg}=1000\alpha,\qquad
 C^*_{\deg}=2(\eta_{\rm ret}+\alpha).
\]

\begin{proposition}[Shift of the base-thousand representation]
\label{prop:ii-desplazamiento-angular}
Let \(\alpha=\sum_{n\geq1}b_n1000^{-n}\), with
\(b_n\in\{0,\ldots,999\}\), be the compatible expansion already generated.
The lifted coordinate \(A_{\deg}=1000\alpha\) satisfies
\[
 A_{\deg}=b_1+\sum_{n\geq1}b_{n+1}1000^{-n}.
\]
Representation preserves the block sequence from the second block onward;
the initial block moves to the integer part.
\end{proposition}
\begin{proof}
The series converges absolutely. Multiplying each term by 1000
and separating the initial term gives the identity. The operation
acts on the constructed expansion, with its blocks and memory;
subsequent circular representation takes its class modulo 360.
\end{proof}

\begin{definition}[Angular-coordinate charts]
The degree chart expresses one complete winding in 360 units.
The radian chart expresses it as \(2\pi\). The coordinate
changes and their inverses are
\[
 x=\frac{\pi}{180}A_{\deg},\qquad
 y=\frac{\pi}{180}C^*_{\deg},\qquad
 A_{\deg}=\frac{180}{\pi}x,\qquad
 C^*_{\deg}=\frac{180}{\pi}y.
\]
On a traversal with retained winding number, the lifted
coordinates are used; the circular quotient expresses their class modulo
the corresponding winding.
\end{definition}

The definition characterizes degrees and radians through their exact relation
to a winding. The register additionally retains orientation, windings, and
transport memory. Representation of its angular class preserves
phase, while the lift distinguishes traversals with the same
final phase. The physical action readout must specify which of
these coordinates enters the integral.

\begin{proposition}[Covariance of the action readout]
For \(h_a=2\pi\hbar_a\), the action of a lifted coordinate
satisfies
\begin{equation}
 \mathcal S_a(\theta)
 =\hbar_a\theta_{\rm rad}
 =h_a\frac{\theta_{\deg}}{360}.
 \label{eq:ii-accion-cartas}
\end{equation}
A positive winding expresses \(h_a\).
\end{proposition}
\begin{proof}
Substituting \(\theta_{\rm rad}=\pi\theta_{\deg}/180\) gives
\(\hbar_a\pi\theta_{\deg}/180
=(2\pi\hbar_a)\theta_{\deg}/360\). In one winding,
\(\theta_{\deg}=360\), yielding \(h_a\).
\end{proof}

The periodicity of a representation is declared separately. If the
transported representation satisfies \(U(2\pi)=-I\), its operator
restoration occurs at \(4\pi\). This property of the lift
coexists with the per-winding normalization in
\eqref{eq:ii-accion-cartas}. Each preserves its object: integrated
quantity, circular phase, and represented state.

\subsection{Covariance of the channels and the incidence functional}

The channels can equivalently be written as
\begin{equation}
 q_\pm=\exp[-(x\pm y)]
 =\exp\!\left[-\frac{\pi}{180}(A_{\deg}\pm C^*_{\deg})\right].
\label{eq:ii-canales}
\end{equation}
\begin{proposition}[Recovery of the even and odd coordinates]
\label{prop:ii-inversion-canales-angulares}
In the lifted real chart, the positive channels uniquely
determine both coordinates:
\[
 A_{\deg}=-\frac{90}{\pi}\log(q_+q_-),\qquad
 C^*_{\deg}=\frac{90}{\pi}\log\frac{q_-}{q_+}.
\]
The interchange \((q_+,q_-)\mapsto(q_-,q_+)\) preserves
\(A_{\deg}\) and reverses \(C^*_{\deg}\).
\end{proposition}
\begin{proof}
The sum and difference of the logarithms in
\eqref{eq:ii-canales} are, respectively,
\(-\pi A_{\deg}/90\) and \(-\pi C^*_{\deg}/90\).
The real logarithm is single-valued on positive channels.
The product is symmetric and the quotient is inverted upon interchange.
\end{proof}
In the twelve-position register, the antipodal involution
\(r\mapsto r+6\pmod {12}\) forms six pairs. Uniform
distribution of the odd coordinate over these pairs is
\(C^*_{\deg}/6\): adding its six contributions recovers
\(C^*_{\deg}\). This sectorial coordinate is distinguished from
the Planck constant \(h\), which belongs to the action line.

The conversion is applied exactly once to each argument. The quantity
\(\alpha\) remains dimensionless and the quantity \(\hbar\)
continues to belong to the action line; their preceding coordinates
enter the map \(\mathcal P\).

In addition to the channels, the Barbero--Immirzi functional contains a
bilinear component in the angular representatives. Its transport is
\begin{equation}
 \frac{\pi}{180}A_{\deg}C^*_{\deg}
 =\frac{180}{\pi}xy.
 \label{eq:ii-barbero-cartas}
\end{equation}
The identity is obtained by replacing both degree representatives
with their radian expressions. Thus the complete formula in
analytic coordinates is
\[
 \gamma_{\HMT}
 =\frac{180}{\pi}xy
 +12\bigl[S_{90}(e^{-x+y})-S_{90}(e^{-x-y})\bigr]
 -\bigl[S_{120}(e^{-x+y})-S_{120}(e^{-x-y})\bigr].
\]
Equivalence of the two charts transports the bilinear component
and exponential functions simultaneously. This transformation
makes explicit the precise role of the degree in the construction and allows
evaluation to be reproduced without ambiguity in angular units.
\par\endgroup
